% POLAR WINDOW: technical paper.
% Build from the repository root:  tectonic docs/paper/paper.tex
% Numbers come from docs/paper/data/ (export_data.py, measure_timing.py) and
% from `make test`, `make test-fe` and `make lint`, run on 2026-09-24.
% IEEEtran asks for the Type 1 families ptm, phv and pcr, which XeTeX cannot
% use; declare them in XeTeX's TU encoding over the TeX Gyre equivalents.
\DeclareFontFamily{TU}{ptm}{}
\DeclareFontShape{TU}{ptm}{m}{n}{<->"[texgyretermes-regular.otf]:mapping=tex-text"}{}
\DeclareFontShape{TU}{ptm}{m}{it}{<->"[texgyretermes-italic.otf]:mapping=tex-text"}{}
\DeclareFontShape{TU}{ptm}{m}{sc}{<->"[texgyretermes-regular.otf]:mapping=tex-text;+smcp"}{}
\DeclareFontShape{TU}{ptm}{m}{sl}{<->ssub * ptm/m/it}{}
\DeclareFontShape{TU}{ptm}{b}{n}{<->"[texgyretermes-bold.otf]:mapping=tex-text"}{}
\DeclareFontShape{TU}{ptm}{b}{it}{<->"[texgyretermes-bolditalic.otf]:mapping=tex-text"}{}
\DeclareFontShape{TU}{ptm}{b}{sc}{<->"[texgyretermes-bold.otf]:mapping=tex-text;+smcp"}{}
\DeclareFontShape{TU}{ptm}{b}{sl}{<->ssub * ptm/b/it}{}
\DeclareFontShape{TU}{ptm}{m}{scit}{<->"[texgyretermes-italic.otf]:mapping=tex-text;+smcp"}{}
\DeclareFontShape{TU}{ptm}{b}{scit}{<->"[texgyretermes-bolditalic.otf]:mapping=tex-text;+smcp"}{}
\DeclareFontShape{TU}{ptm}{bx}{scit}{<->ssub * ptm/b/scit}{}
\DeclareFontShape{TU}{ptm}{bx}{n}{<->ssub * ptm/b/n}{}
\DeclareFontShape{TU}{ptm}{bx}{it}{<->ssub * ptm/b/it}{}
\DeclareFontShape{TU}{ptm}{bx}{sc}{<->ssub * ptm/b/sc}{}
\DeclareFontShape{TU}{ptm}{bx}{sl}{<->ssub * ptm/b/it}{}
\DeclareFontFamily{TU}{phv}{}
\DeclareFontShape{TU}{phv}{m}{n}{<->"[texgyreheros-regular.otf]:mapping=tex-text"}{}
\DeclareFontShape{TU}{phv}{m}{it}{<->"[texgyreheros-italic.otf]:mapping=tex-text"}{}
\DeclareFontShape{TU}{phv}{b}{n}{<->"[texgyreheros-bold.otf]:mapping=tex-text"}{}
\DeclareFontShape{TU}{phv}{b}{it}{<->"[texgyreheros-bolditalic.otf]:mapping=tex-text"}{}
\DeclareFontShape{TU}{phv}{bx}{n}{<->ssub * phv/b/n}{}
\DeclareFontShape{TU}{phv}{bx}{it}{<->ssub * phv/b/it}{}
\DeclareFontFamily{TU}{pcr}{\hyphenchar\font=-1}
\DeclareFontShape{TU}{pcr}{m}{n}{<->"[texgyrecursor-regular.otf]"}{}
\DeclareFontShape{TU}{pcr}{m}{it}{<->"[texgyrecursor-italic.otf]"}{}
\DeclareFontShape{TU}{pcr}{b}{n}{<->"[texgyrecursor-bold.otf]"}{}
\DeclareFontShape{TU}{pcr}{b}{it}{<->"[texgyrecursor-bolditalic.otf]"}{}
\DeclareFontShape{TU}{pcr}{bx}{n}{<->ssub * pcr/b/n}{}
\DeclareFontShape{TU}{pcr}{bx}{it}{<->ssub * pcr/b/it}{}
\DeclareFontShape{TU}{pcr}{m}{sc}{<->ssub * pcr/m/n}{}
\DeclareFontShape{TU}{pcr}{b}{sc}{<->ssub * pcr/b/n}{}
\DeclareFontShape{TU}{pcr}{bx}{sc}{<->ssub * pcr/b/n}{}
\documentclass[conference]{IEEEtran}

\usepackage{amsmath}
\usepackage{amssymb}
\usepackage{newtxmath}
\usepackage{booktabs}
\usepackage{graphicx}
\usepackage[per-mode=symbol]{siunitx}
\sisetup{
  uncertainty-mode = separate,
  group-minimum-digits = 4,
  range-units = single,
}
% algorithm2e.sty carries Latin-1 bytes; read it as Latin-1 so XeTeX does not
% warn about each one, then return to UTF-8 for everything that follows.
\XeTeXdefaultencoding "iso-8859-1"
\usepackage[ruled,linesnumbered,vlined]{algorithm2e}
\XeTeXdefaultencoding "utf-8"
\usepackage{pgfplots}
\pgfplotsset{compat=1.18}
\usepackage{microtype}
\usepackage[backend=biber,style=ieee]{biblatex}
\addbibresource{references.bib}
% Let long URLs break anywhere rather than leave gaps in a narrow column.
\setcounter{biburlnumpenalty}{9000}
\setcounter{biburlucpenalty}{9000}
\setcounter{biburllcpenalty}{9000}
\AtBeginBibliography{\raggedright}
\usepackage{cleveref}
\crefname{algocf}{Algorithm}{Algorithms}
\Crefname{algocf}{Algorithm}{Algorithms}

% An external figure, or a framed notice naming the file while it is pending.
\newcommand{\extfig}[2]{%
  \IfFileExists{figures/#1}%
    {\includegraphics[width=\columnwidth]{figures/#1}}%
    {\fbox{\parbox[c][#2][c]{0.94\columnwidth}{\centering\footnotesize
      Figure pending: \texttt{figures/#1}}}}}

\pgfplotsset{
  pw/.style={
    width=\columnwidth, height=0.62\columnwidth,
    tick label style={font=\footnotesize},
    label style={font=\footnotesize},
    legend style={font=\footnotesize, draw=none, fill=none},
    axis line style={black!70}, tick style={black!70},
  },
}

\begin{document}

\title{Terrain, Not Latitude: Sunlight and Earth Contact at South-Polar Landing
Sites, With the Uncertainty Each Number Carries}

\author{\IEEEauthorblockN{Md Ahbab}
\IEEEauthorblockA{POLAR WINDOW\\NASA Space Apps Challenge 2026,
CLPS Lunar Mission Browser}}

\maketitle

\begin{abstract}
Commercial Lunar Payload Services (CLPS) deliveries are converging on the lunar
south pole, where the Sun circles within a few degrees of the horizon and
terrain, not latitude, decides whether a lander has power and a line of sight to
Earth. We describe the geometry engine of POLAR WINDOW, a CLPS mission browser
that reports four numbers for any site and window: hours of sunlight, hours of
Earth contact, the longest unbroken dark gap and the longest unbroken radio
blackout, never combined into a score. Each skyline is ray-marched once from
Lunar Orbiter Laser Altimeter terrain; everything after is a lookup and one
vectorised SPICE call per body. The contribution is the uncertainty model:
per-bearing elevation error, a near-field noise floor, grazing crossings reported
as untimeable, and window-censored gaps reported as lower bounds. We also measure
what the tolerance omits: a tuning constant moves a summit's sunlight from
\qty{574.399}{\hour} to \qty{482.33}{\hour}.
\end{abstract}

\section{Problem}

At Connecting Ridge (\ang{-89.438} latitude) the Sun never rises: across the
\qty{708.75}{\hour} window from 2027-03-14 used throughout, its elevation stays
between \ang{-1.607} and \ang{-0.7031} (\Cref{fig:sunskyline}). A sun calculator
with a level horizon reports no sunlight. Yet the site is sunlit for
\qty{81.04}{\percent} of the window, because it stands \qty{1957.5}{\metre} above
the \qty{1737.4}{\kilo\metre} reference sphere and looks \emph{down} at its
skyline, whose mean over \num{1440} bearings is \ang{-2.3681}. At the pole the
horizon decides~\cite{mazarico2011illumination,bussey2010kaguya,speyerer2013persistently,glaser2014illumination}.

The second trap is the coordinate. A crater's published coordinate is its
centre, which is its floor: the \qty{80}{\metre} LOLA model returns
\qty{-2767.3}{\metre}, \qty{-3315.9}{\metre} and \qty{-3931.7}{\metre} at the
anchors for Shackleton, Haworth and Shoemaker. The \num{15} candidate sites are
therefore resolved from terrain: the most open-skied post, sloped at most
\ang{12}, in the top elevation decile around each anchor. Eight of the \num{20}
records in the project's CLPS dataset target the south pole.

\begin{figure}[t]
\centering
\begin{tikzpicture}
\begin{axis}[pw,
  xlabel={Hours since 2027-03-14 00:00 UTC},
  ylabel={Elevation (\si{\degree})},
  xmin=0, xmax=708.75, ymin=-4.6, ymax=0.4,
  xtick={0,100,200,300,400,500,600,700},
  legend style={at={(0.5,1.03)}, anchor=south, legend columns=2,
    /tikz/every even column/.append style={column sep=8pt}},
  legend cell align=left]
\addplot[black!45, thin, densely dotted, forget plot] coordinates {(0,0) (708.75,0)};
\addplot[black!55, thick, dashed] table[x=t_h, y=sun_skyline_deg] {sun_vs_skyline.dat};
\addlegendentry{Skyline at the Sun's azimuth}
\addplot[black, thick] table[x=t_h, y=sun_el_deg] {sun_vs_skyline.dat};
\addlegendentry{Sun elevation}
\end{axis}
\end{tikzpicture}
\caption{At Connecting Ridge the Sun stays below level (dotted) for the whole
window, yet it clears the site's own skyline for \qty{81.04}{\percent} of it,
because that skyline lies below level on every bearing. Data:
\texttt{sun\_vs\_skyline.dat}, \num{2836} samples at \qty{900}{\second}.}
\label{fig:sunskyline}
\end{figure}

\section{Method}

\paragraph{Horizon mask} \Cref{alg:raymarch} marches \num{1440} rays at
\ang{0.25} to \qty{50}{\kilo\metre}, \num{555} samples each, starting two posts
out, beyond the site's own pixel. Each elevation angle
subtracts the curvature drop,
\begin{equation}
\theta_k = \arctan\frac{h(d_k) - h_0 - d_k^2/2R}{d_k},\qquad
R = \qty{1737.4}{\kilo\metre},
\label{eq:theta}
\end{equation}
which is \qty{115.11}{\metre} at \qty{20}{\kilo\metre}, and is never clamped at
zero (\Cref{fig:raymarch}). Terrain is LDEM\_80S\_20M within
\qty{20}{\kilo\metre} and LDEM\_80S\_80M beyond~\cite{lola_gdr_polar,smith2010lola}.

\paragraph{Visibility} One SPICE~\cite{acton1996naif,acton2018spice}
\texttt{spkpos} call per body, through SpiceyPy~\cite{annex2020spiceypy}, gives
Sun and Earth from DE440~\cite{park2021de440} in MOON\_ME~\cite{naif_moon_fk},
projected onto the site's east, north and up unit vectors, which stay exact at
the pole. The margin $m = e - H(a)$ subtracts the mask at the body's azimuth
from its elevation. Discs are finite: the angular radius $\rho$ spans
\ang{0.2660} to \ang{0.2683} for the Sun and \ang{0.9025} to \ang{0.9963} for
Earth, so a sample is lit if $m > \rho$, dark if $m < -\rho$ and otherwise
penumbral (\Cref{fig:states}), as a median \qty{8.44}{\percent} of samples are.

\paragraph{Intervals and metrics} Sign changes of $m$ are bracketed at
\qty{900}{\second} and bisected \num{40} times; half-open runs partition the
window. The four numbers are two totals and two longest runs, never combined.

\paragraph{Why precompute} The mask depends only on the site and costs a median
\qty{128.653}{\milli\second}; a cold \texttt{/metrics} request costs
\qty{10.162}{\milli\second} (\Cref{tab:timing}; \Cref{fig:timing}). The mask is
built once per site and engine version, and a request is a lookup plus one
ephemeris call.

\begin{algorithm}[t]
\caption{Horizon mask for one site}
\label{alg:raymarch}
\DontPrintSemicolon
\KwIn{height sampler $h(x,y)$, vertical-error sampler $\sigma(x,y)$, site
height $h_0$, post spacing $p$}
\KwOut{per bearing $a$: skyline $\theta^*(a)$, range $r^*(a)$, tolerance $\delta(a)$}
$d_1 \gets 2p$;\ $d_{k+1} \gets d_k + \max(\qty{5}{\metre},\, 0.01\,d_k)$ while
$d_{k+1} \le \qty{50}{\kilo\metre}$\;
\For{$a \in \{0, \ang{0.25}, \dots, \ang{359.75}\}$}{
  $(x_k, y_k) \gets d_k(\sin a, \cos a)$ for all $k$\;
  \lIf{no $h(x_k,y_k)$ is defined}{$\theta^*(a) \gets$ smooth-sphere dip; \textbf{continue}}
  $\theta_k \gets$ \Cref{eq:theta};\quad $\phi_k \gets \arctan(\sigma_k / d_k)$\;
  $E \gets \{k : \theta_k > \phi_k \ \lor\ \phi_k < \ang{0.25}\}$\;
  $k^* \gets \arg\max_{k \in E} \theta_k$\;
  $\theta^*(a) \gets \theta_{k^*}$;\ $r^*(a) \gets d_{k^*}$;\
  $\delta(a) \gets \arctan(\sigma_{k^*}/d_{k^*})$\;
}
\end{algorithm}

\begin{figure}[t]
\centering
\extfig{raymarch-geometry.pdf}{0.55\columnwidth}
\caption{Each sample's elevation angle is measured after the surface has curved
away by $d^2/2R$, so a site above its surroundings has a skyline below level,
and the mask keeps that negative angle.}
\label{fig:raymarch}
\end{figure}

\begin{figure}[t]
\centering
\extfig{state-decision.pdf}{0.5\columnwidth}
\caption{A sample is lit or dark only when the whole disc is clear of, or hidden
by, the skyline; the band in between is reported as penumbral, not rounded to
either side.}
\label{fig:states}
\end{figure}

\begin{table}[t]
\centering
\caption{Precompute versus request cost, Connecting Ridge, reference window at
\qty{900}{\second}. Apple M3 (Mac15,12), macOS 27.0, Python 3.12.14.
Values as written to \texttt{data/timing.csv}.}
\label{tab:timing}
\footnotesize
\begin{tabular}{@{}l S[table-format=3.3] S[table-format=3.3] S[table-format=2.0]@{}}
\toprule
Stage & {Median (\si{\milli\second})} & {p95 (\si{\milli\second})} & {$N$} \\
\midrule
\multicolumn{4}{@{}l}{\emph{Once per site (worker)}} \\
Ray-march, \num{1440} bearings & 128.653 & 130.154 & 75 \\
Load \qty{20}{\metre} window & 22.223 & 25.435 & 45 \\
Load \qty{80}{\metre} product & 399.12 & 420.001 & 3 \\
Fetch \qty{20}{\metre} window from PDS & \multicolumn{3}{c}{not measured} \\
\midrule
\multicolumn{4}{@{}l}{\emph{Per request (API)}} \\
Ephemeris, two bodies & 6.552 & 6.704 & 50 \\
Sun track and crossings & 1.938 & 2.067 & 50 \\
Earth track and crossings & 0.694 & 0.707 & 50 \\
Four numbers and autonomy & 0.008 & 0.010 & 50 \\
HTTP \texttt{/metrics}, cache miss & 10.162 & 11.249 & 30 \\
HTTP \texttt{/metrics}, cache hit & 0.782 & 0.884 & 30 \\
HTTP \texttt{/timeline}, cache miss & 14.369 & 14.930 & 30 \\
\bottomrule
\end{tabular}
\end{table}

\begin{figure}[t]
\centering
\extfig{timing-breakdown.pdf}{0.5\columnwidth}
\caption{The mask costs \num{12.66} times a complete cold request and depends
only on the site, so it is built once; every request is a lookup and an
ephemeris call.}
\label{fig:timing}
\end{figure}

\section{Uncertainty}

\paragraph{Per-tier error} Each bearing carries the vertical error $\sigma$ of
the tier that set its skyline, \qty{2}{\metre} on the \qty{20}{\metre} tier or
\qty{4}{\metre} on the \qty{80}{\metre} tier, as $\delta(a) = \arctan(\sigma/r^*)$.
At Connecting Ridge \num{746} of \num{1440} bearings are set on the coarse tier,
so the site reports \qty{4}{\metre}, not the finest tier's figure.

\paragraph{Near-field noise floor} A \qty{2}{\metre} error at \qty{40}{\metre} is
\ang{2.8624} of apparent skyline (\Cref{fig:noisefloor}). A sample may set the
skyline only if it clears that floor, or if the floor is below \ang{0.25}, which
holds beyond \qty{458.4}{\metre} for \qty{2}{\metre}.

\paragraph{Grazing crossings} An edge's timing tolerance is
$\delta t = \delta(a)/\lvert \mathrm{d}m/\mathrm{d}t \rvert$ (\Cref{fig:polar}).
Below a crossing rate of \qty{0.0025}{\degree\per\hour}, or with $\delta t$ above
\qty{24}{\hour}, the crossing time
is not knowable from these data, and the edge is reported as \emph{grazing},
without a tolerance, rather than with a large one. Faustini rim's longest dark
gap and blackout both graze (\Cref{tab:three}; \Cref{fig:four}).

\paragraph{Censoring} A run touching the window edge was ended by the window,
so its duration is a lower bound. Five of the \num{15} sites have a censored
longest night. A bound claims $[v, \infty)$; an estimate claims $[v-u, v+u]$. A
site beats another only when its interval clears the other's, so on longest
night Connecting Ridge (\qty{121.57 \pm 17.09}{\hour}) ranks 1--2, de Gerlache
rim ($\geqslant$\,\qty{112.17}{\hour}) 1--3 and Malapert massif
($\geqslant$\,\qty{287.76}{\hour}) 2--3. Reading the bound as an estimate would
make de Gerlache rim the best site and size its battery for a night shorter than
the real one.

\paragraph{What the tolerance omits} On \num{424} of Connecting Ridge's bearings
the ground falls away from the summit, so the skyline is set at
\qty{460}{\metre}, the first sample the floor trusts: the rule, not a ridge,
fixes it. Both edges of the longest night lie there, hence
\qty{\pm 17.09}{\hour}. Moving the threshold to \ang{0.125} or \ang{0.5} changes
sunlight to \qty{611.823}{\hour} or \qty{482.33}{\hour} (\Cref{tab:sensitivity}),
far outside that tolerance.

\begin{figure}[t]
\centering
\begin{tikzpicture}
\begin{loglogaxis}[pw,
  xlabel={Range from site (\si{\metre})},
  ylabel={Angular noise floor (\si{\degree})},
  xmin=20, xmax=5000, ymin=0.02, ymax=12,
  legend pos=north east, legend cell align=left]
\addplot[black, thick] table[x=range_m, y=floor2] {noise_floor.dat};
\addlegendentry{$\sigma = \qty{2}{\metre}$}
\addplot[black!55, thick, dashed] table[x=range_m, y=floor4] {noise_floor.dat};
\addlegendentry{$\sigma = \qty{4}{\metre}$}
\addplot[black!45, densely dotted, forget plot] coordinates {(20,0.25) (5000,0.25)};
\addplot[only marks, mark=*, mark size=1.5pt, black, forget plot] coordinates {(40,2.86241)};
\draw[black!60, thin] (axis cs:40,2.86241) -- (axis cs:40,0.54);
\node[anchor=west, font=\footnotesize, fill=white, inner sep=1pt] at (axis cs:26,0.45) {\ang{2.8624} at \qty{40}{\metre}};
\node[anchor=south east, font=\footnotesize] at (axis cs:4800,0.25) {\ang{0.25}};
\end{loglogaxis}
\end{tikzpicture}
\caption{Close to the site, elevation error alone produces skylines of several
degrees. Samples are trusted unconditionally only where the floor
$\arctan(\sigma/r)$ falls below \ang{0.25} (dotted): beyond \qty{458.4}{\metre}
for \qty{2}{\metre} and \qty{916.7}{\metre} for \qty{4}{\metre}. Data:
\texttt{noise\_floor.dat}.}
\label{fig:noisefloor}
\end{figure}

\begin{figure}[t]
\centering
\extfig{horizon-polar-connecting-ridge.pdf}{0.9\columnwidth}
\caption{Connecting Ridge's skyline lies below level on every bearing. The
crossings that bound its longest night fall where the skyline tolerance is
largest, so they carry the widest timing bands.}
\label{fig:polar}
\end{figure}

\begin{table}[t]
\centering
\caption{The four numbers for three contrasting sites over the reference window,
in hours, rounded to two decimals from \texttt{data/sites\_four\_numbers.csv}.
$\geqslant$ marks a lower bound; \emph{grazing} marks an edge that cannot be timed.}
\label{tab:three}
\footnotesize
\setlength{\tabcolsep}{4pt}
\newcommand{\hd}[1]{\begin{tabular}[b]{@{}l@{}}#1\end{tabular}}
\begin{tabular}{@{}llll@{}}
\toprule
Metric (\si{\hour}) & \hd{Connecting\\Ridge} & \hd{Malapert\\massif} & \hd{Faustini\\rim} \\
\midrule
Sunlight & \num{574.40 \pm 17.09} & \num{394.16 \pm 7.25} & \num{103.07 \pm 24.00} \\
Earth contact & \num{430.28 \pm 5.26} & \num{570.61 \pm 7.63} & \num{278.57 \pm 33.94} \\
Longest dark & \num{121.57 \pm 17.09} & $\geqslant$\,\num{287.76} & \num{605.68}, grazing \\
Longest blackout & \num{278.47 \pm 5.26} & \num{138.14 \pm 7.63} & \num{430.18}, grazing \\
\bottomrule
\end{tabular}
\end{table}

\begin{figure}[t]
\centering
\extfig{four-numbers-three-sites.pdf}{0.6\columnwidth}
\caption{A measured value carries an error bar, a censored value is drawn as a
bound with an open end, and a grazing value has no bar at all: three different
claims that a single point would hide.}
\label{fig:four}
\end{figure}

\begin{table}[t]
\centering
\caption{Connecting Ridge under three values of the near-field threshold, as
written to \texttt{data/sensitivity.json}. The engine uses \ang{0.25}.}
\label{tab:sensitivity}
\footnotesize
\setlength{\tabcolsep}{4pt}
\begin{tabular}{@{}S[table-format=1.3] S[table-format=3.1] S[table-format=-1.4]
  S[table-format=3.3] S[table-format=3.3]@{}}
\toprule
{Threshold} & {Trusted from} & {Mean skyline} & {Sunlight} & {Longest dark} \\
{(\si{\degree})} & {(\si{\metre})} & {(\si{\degree})} & {(\si{\hour})} & {(\si{\hour})} \\
\midrule
0.125 & 916.7 & -2.6566 & 611.823 & 84.145 \\
0.25 & 458.4 & -2.3681 & 574.399 & 121.569 \\
0.5 & 229.2 & -1.7756 & 482.33 & 206.515 \\
\bottomrule
\end{tabular}
\end{table}

\section{Validation}

\paragraph{Golden fixtures} Three sites chosen for different physics are rebuilt
from terrain on disk, not compared with a stored mask, to \ang{0.05}
on the skyline, \qty{72}{\second} on arc edges and the larger of
\qty{360}{\second} and \qty{0.1}{\percent} on the four numbers. Connecting
Ridge's rebuilt mask is identical to the stored one.

\paragraph{Properties} Interval code knows no astronomy, so
Hypothesis~\cite{maciver2019hypothesis} draws up to \num{300} boolean series per
property and checks partition, alternation, contiguity, inversion symmetry and
monotonicity, with degenerate shapes pinned.

\paragraph{Adapter} The browser's adapter tests assert its hours against the
seconds the backend golden pins.

\paragraph{Engine stamp} The engine version, hz-2/vis-2, is in every cache key,
stored mask and response, so geometry from another engine is a readable 409,
not a stale answer. \Cref{tab:tests} lists the results; the one failure is the
OpenAPI snapshot, stale against uncommitted schema changes.

\begin{table}[t]
\centering
\caption{What \texttt{make test}, \texttt{make test-fe} and \texttt{make lint}
printed on 2026-09-24, with the backend suites also run per directory.}
\label{tab:tests}
\footnotesize
\begin{tabular}{@{}l l S[table-format=3.0] S[table-format=1.0]@{}}
\toprule
Command & Suite & {Passed} & {Failed} \\
\midrule
\texttt{make test} & unit & 159 & 0 \\
 & integration & 56 & 1 \\
 & golden & 16 & 0 \\
 & total, \qty{6.41}{\second} & 231 & 1 \\
\texttt{make test-fe} & adapter, keys, exports & 41 & 0 \\
\midrule
\texttt{make lint} & \multicolumn{3}{l}{ruff: no findings, \num{80} files formatted} \\
 & \multicolumn{3}{l}{mypy strict: no issues in \num{80} files} \\
 & \multicolumn{3}{l}{import contracts: \num{8} kept, \num{0} broken} \\
\bottomrule
\end{tabular}
\end{table}

\section{Limitations}

The \qty{120}{\hour} survival limit is a browser placeholder for a thermal model
the system does not have; a per-lander thermal and battery model would replace
it. Elevation error is constant per tier, not per pixel; the LOLA count products
(LDEC\_80S\_20M and LDEC\_80S\_80M~\cite{lola_gdr_polar}) would localise it. At
summits the near-field threshold outweighs the reported tolerance; a
lander-height term and \qty{5}{\metre} site maps~\cite{barker2021improved} would
bound it. Illumination and shadow tiles use a five-ring horizon proxy and are
labelled so; per-pixel masks on the \qty{80}{\metre} grid mean \num{57760000}
ray-marches, \qty{86.0}{\day} single-threaded. No screen-reader user has tested
the arc announcements; only such users can. Fetching a \qty{20}{\metre} window
from PDS was not timed; timing \texttt{make seed} on an empty volume would.

\section{Conclusion}

POLAR WINDOW answers a CLPS planner with four numbers per site and window, each
carrying a tolerance, the word \emph{grazing}, or a lower bound, and none folded
into a score. Geometry costs \qty{128.653}{\milli\second} once per site; a cold
request then takes \qty{10.162}{\milli\second}. The uncertainty model separates
crossings the data can time from those they cannot, and bounds from estimates,
which changes how sites rank. Its own sensitivity test names the next step: at
summits the near-field rule moves sunlight by more than the reported tolerance,
and that must be bounded before these numbers choose a landing site.

\printbibliography

\end{document}
